\section{The Expressiveness Threshold}
\label{sec:model}

\paragraph{Setup.}
Two agents play a repeated game with per-round action sets $A_i$ and a message space
$M$ available before each action. Classical cheap-talk analysis of the Prisoner's
Dilemma treats $M$ as strategically inert: announcements are unverifiable and add no
equilibrium~\cite{crawford1982strategic}. Folk-theorem
reasoning~\cite{friedman1971noncooperative,fudenberg1986folk} says cooperative paths
become sustainable when deviation triggers a punishment continuation. Both accounts
make predictions about $M$ that the data can separate.

\paragraph{Two classes of message.}
Distinguish, within $M$, messages by the \emph{grammatical person of the commissive}
they can express:
\begin{description}[leftmargin=2.4em,style=nextline,itemsep=1pt,topsep=2pt]
\item[\textnormal{\textbf{Own intention}}] a first-person-singular commissive naming
only the sender's own next action: ``I will cooperate.'' It conveys a prediction about
one agent.
\item[\textnormal{\textbf{Joint proposal}}] a first-person-plural commissive naming an
action profile for \emph{both} parties: ``let's both hold at $12.00$.'' It nominates a
point in the joint action space and invites acceptance.
\end{description}

\paragraph{The claim.}
Coordination requires $M$ to admit a joint proposal. The mechanism is equilibrium
\emph{selection}, not enforcement: a repeated game with multiple equilibria leaves the
agents facing a selection problem that an own-intention announcement cannot solve,
because it does not name a profile. A joint proposal names one, and acceptance is
common knowledge once both have seen it.

We state this as a property of the message space rather than of its size. Let
$a^\ast \in A_1 \times A_2$ be the cooperative profile. Write $M$ as \emph{profile-naming}
if some $m \in M$ denotes $a^\ast$ as a profile for both players, and
\emph{intention-naming} if every $m \in M$ denotes only a component $a_i$.

\begin{quote}
\textbf{Predicate hypothesis (pre-registered, and refuted in \S\ref{sec:p1}).}
Coordination on $a^\ast$ requires $M$ to be profile-naming. Under an intention-naming
$M$, an announcement is a unilateral prediction: it does not make $a^\ast$ mutually
salient, so play remains at the risk-dominant equilibrium regardless of $|M|$.
Capacity beyond the shortest profile-naming message is inert.
\end{quote}

This is a claim about \emph{selection among equilibria}, not about \emph{sustaining}
one, and it makes no appeal to punishment. It is closer to focal-point
reasoning~\cite{schelling1960strategy} than to the folk theorem, and it predicts the
opposite of the folk-theorem gloss for exactly the cases \S\ref{sec:reanalysis} tests.

\paragraph{We state it because we tested it, and it failed.}
P1 (\S\ref{sec:p1}) manipulated exactly this property with length, authorship and
per-round availability held fixed. Its pre-registered primary contrast came out as
predicted, and its control arm refuted the hypothesis anyway: a fixed menu whose entire
vocabulary is ``I intend to cooperate''/``I intend to defect''---intention-naming by
construction---produced the \emph{highest} coordination of any arm. We report the
predicate as refuted rather than quietly softening it, and we keep it in the paper
because what replaces it is only interpretable against it.

\paragraph{What survives: the channel is a decision-shaping device.}
Across the restricted arms of P1 the messages are near-templated, so what an agent
\emph{announced} can be read directly. Announcements track actions almost
one-to-one, and agents keep them $86$--$97\%$ of the time. The arms do not differ in
what a message can \emph{encode}; they differ in what the framing leads the agent to
\emph{decide}. A joint-proposal frame yields $87\%$ cooperative announcements; an
unscaffolded first-person-singular frame yields $58\%$, because an agent forbidden to
refer to its counterpart reasons about itself and commits accordingly; a curated menu
yields $94.5\%$, not through expressiveness but through \emph{choice architecture}---a
pre-written cooperative option placed in front of the agent.

The revised account is therefore not about channel capacity or channel structure at
all:

\begin{quote}
\textbf{The channel acts on the agent's decision, not on the peer's information.}
What an interface makes easy to say changes what the agent resolves to do, and the
agent then does it. Two very different interfaces---an open proposal frame and a
two-item menu with a salient cooperative option---converge because both make
cooperation the natural thing to utter.
\end{quote}

This is weaker than the predicate we set out to establish, and it is correlational
where the predicate was structural: P1 manipulates the frame and observes the
announcement shift, but announcement and action are chosen by the same forward pass, so
we cannot separate ``the frame changed the decision'' from ``the frame changed the
report of a decision already made''. We flag the experiment that would separate
them---forcing the action choice before the message is composed---as untested.

\paragraph{Three predictions that separate this from the alternatives.}
\begin{enumerate}[leftmargin=1.4em,itemsep=2pt,topsep=2pt]
\item \textbf{Against bandwidth.} Capacity should not matter once the floor is
cleared: a very short proposal should beat a much longer own-intention message.
\item \textbf{Against enforcement.} Contingent punishment rules should be neither
necessary nor sufficient---coordination should occur without them, and their presence
should not predict it.
\item \textbf{Threshold, not dose.} Among channels that clear the floor, variation in
how much or how well agents talk should stop predicting the outcome.
\end{enumerate}
Sections~\ref{sec:reanalysis}--\ref{sec:p3} test all three.

% ======================================================================
